\def\writeANS{}
\begin{loigiaibth}{12}
  Theo phương trình nhiệt hóa học ta có: Cứ đốt cháy $1$ mol CO thì nhiệt lượng tỏa ra là $851{,}15\, kJ$. \par (a) Ta có: $n_{CO} = \dfrac {2{,}479}{24{,}79} = 0{,}1\, mol \Rightarrow $ Nhiệt lượng tỏa ra là: $0{,}1.851{,}5 = 85{,}15\, kJ$ \par (b) Ta có: $n_{CO} = \dfrac {1277{,}25}{851{,}5} = 1{,}5\, mol \Rightarrow $ Thể tích CO đã dùng ở đkc là $V_{CO} = 1{,}5.24{,}79 = 37{,}185(L)$.
\end{loigiaibth}
\def\writeANS{}
\begin{loigiaibth}{13}
  Từ phản ứng ta thấy \par Cứ $1$ mol $N_2$ phản ứng hết tạo ra $2$ mol $NH_3$ và tỏa ra $99{,}22$ kJ \par Cần $0{,}5$ mol $N_2$ phản ứng hết tạo ra $1$ mol $NH_3$ và tỏa ra x kJ \par $\Rightarrow x = \dfrac {99{,}22}{2} = 49{,}61kJ$ \par $\Rightarrow \Delta _f H_{298}^\circ = -49{,}61kJ\, mol^{-1}$ (enthalpy có giá trị âm vì đây là phản ứng tỏa nhiệt)
\end{loigiaibth}
\def\writeANS{}
\begin{loigiaibth}{14}
  \begin {enumerate} \item [(a)] $\Delta _r H_{298}^\circ (1) = 2.(-46) - 1.0 - 3.0 = -92\, kJ$ \par $\Delta _r H_{298}^\circ (2) = (-46) - \dfrac {1}{2}.0 - \dfrac {3}{2}.0 = -46\, kJ$ \par Phản ứng tỏa nhiệt và $\Delta _r H_{298}^\circ (1) = 2. \Delta _r H_{298}^\circ (2)$ \item [(b)] Khi tổng hợp $1$ tấn $NH_3$ thì nhiệt lượng tỏa ra $= \dfrac {46.10^6}{17} = 2{,}7.10^6 (kJ)$ \par Tính theo $2$ phương trình trên thì đều ra kết quả giống nhau. \end {enumerate}
\end{loigiaibth}
\def\writeANS{}
\begin{loigiaibth}{15}
  \begin {enumerate} \item [(a)] Nhân phương trình của phản ứng với 3: $\Delta _r H_{298}^\circ = -285{,}66 x 3 = -856{,}98\, kJ$ \item [(b)] Chia phương trình của phản ứng với 2: $\Delta _r H_{298}^\circ = -285{,}66 : 2 = -142{,}83\, kJ$ \item [(c)] Đảo chiều của phản ứng: $\Delta _r H_{298}^\circ = +285{,}66\, kJ$ \end {enumerate}
\end{loigiaibth}
\def\writeANS{}
\begin{loigiaibth}{16}
  \begin {enumerate} \item [(a)] $H_2(g) + \dfrac {1}{2}O_2(g) \longrightarrow H_2O(g)\,\, \Delta _r H_{298}^\circ = -214{,}6\, kJ$ \item [(b)] $H_2(g) + \dfrac {1}{2}O_2(g) \longrightarrow H_2O(l)\,\, \Delta _r H_{298}^\circ = -285{,}49\, kJ$ \item [(c)] Số mol ammonia $= \dfrac {5}{34}$ mol. $1$ mol ammonia tỏa ra $22{,}99 : \dfrac {5}{34} = 156{,}332$ kJ nhiệt. \par $\dfrac {3}{2}H_2(g) + \dfrac {1}{2}N_2(g) \longrightarrow NH_3(g)\,\, \Delta _r H_{298}^\circ = -156{,}332\, kJ$ \item [(d)] $n_{CaO} = 0{,}2$ mol $\Rightarrow $ Để tạo thành $1$ mol CaO cần cung cấp $6{,}94.5 = 34{,}7$ kcal. \par $CaCO_3(s) \xrightarrow {t^\circ } CO_2(g) + CaO(s)\,\, \Delta _r H_{298}^\circ = +34{,}7$ kcal \end {enumerate}
\end{loigiaibth}
\def\writeANS{}
\begin{loigiaibth}{17}
  Tính khối lượng lactic acid tạo ra từ quá trình chạy bộ. Năng lượng của sự chuyển hoá glucose thành lactic acid trong quá trình chạy bộ chiếm $2\% x 300$ kcal $= 6$ kcal $= 6 000$ cal $\Leftrightarrow 25 104 J = 25{,}104\, kJ$. \par \begin {tabular}{ccc} $C_6H_{12}O_6$ & $\longrightarrow 2C_3H_6O_3$ & $\Delta _r H_{298}^\circ = -150\, \text {kJ}$ \\ $0{,}335$ mol & & $-25{,}104$ kJ \end {tabular} \par Khối lượng lactic acid được tạo ra trong quá trình chuyển hóa: $0{,}335 x 90 = 30{,}15$ g.
\end{loigiaibth}
\def\writeANS{}
\begin{loigiaibth}{18}
  Ta có: (1) – (2) = (3) $\Rightarrow \Delta _r H_{298}^\circ (3) = \Delta _r H_{298}^\circ (1) - \Delta _r H_{298}^\circ (2) = -393{,}5 - 2{,}87 = -396{,}37\, kJ$
\end{loigiaibth}
\def\writeANS{}
\begin{loigiaibth}{19}
  Ta có: (3) $H_2(g) + \dfrac {1}{2}O_2(g) \longrightarrow H_2O(l)\,\, \Delta _r H_{298}^\circ (3) = -286kJ$ \par (4) $Mg(s) + \dfrac {1}{2}O_2(g) \longrightarrow MgO(s)\,\, \Delta _r H_{298}^\circ (4) = ?\, kJ$. \par Ta có: (4) = (1) – (2) + (3) $\Rightarrow \Delta _r H_{298}^\circ (4) = -467 - (-151) + (-286) = -602\, kJ$ \par $\Rightarrow $ Enthalpy tạo thành chuẩn của MgO(s) là $\Delta _f H_{298}^\circ $ (MgO, s) $= -602$ kJ/mol
\end{loigiaibth}
\def\writeANS{}
\begin{loigiaibth}{20}
  \begin {enumerate} \item [(a)] Quá trình tan chảy của nước đá là quá trình thu nhiệt vì có biến thiên enthalpy dương. \item [(b)] Viên đá lại tan chảy dần vì nó lấy nhiệt từ nước lỏng( là môi trường xung quanh) \item [(c)] Nước lỏng nhường nhiệt cho viên nước đá, sự mất nhiệt làm cho nước lỏng lạnh đi. \item [(d)] Nhiệt lượng mà $500$ gam nước lỏng từ ở $20^\circ C$ xuống $0^\circ C$ tỏa ra là: \par $\dfrac {500}{18}.75{,}4.(20-0) = 41888{,}9 (J) = 41{,}8889 (kJ)$ \par Mỗi viên đá tan chảy hấp thụ một nhiệt lượng là $6{,}020$ kJ \par $\Rightarrow $ Số viên nước đá tối thiểu cần là $\dfrac {41{,}8889}{6{,}02} = 6{,}96 \approx 7$ viên. \item [(e)] Nhiệt lượng mà $120gam$ nước lỏng ở $45^\circ C$ xuống $0^\circ C$ tỏa ra là: \par $\dfrac {120}{18}.75{,}4.(45-0) = 22620 (J) \Rightarrow $ Lượng nước đá cần dùng là: $(22620 : 6020) x 18 = 67{,}63 (g)$ \par Vậy dùng $150$ gam nước đá là dư \end {enumerate}
\end{loigiaibth}
\def\writeANS{}
\begin{loigiaibth}{21}
  Lượng nhiệt cần dùng để đun nóng $500$ gam nước từ $20^\circ C$ lên $90^\circ C$ là \par $\dfrac {500}{18}.75{,}4.(363-293) = 146611{,}1(J)$ \par Vậy lượng than đá cần dùng là: $\dfrac {146611{,}1}{23.10^3} = 6{,}37 (g)$.
\end{loigiaibth}
\def\writeANS{}
\begin{loigiaibth}{22}
  Để thu được $1000$ kg $CaO$ thì nhiệt lượng cần cung cấp (với hiệu suất là $60\%$) \\ $\dfrac {1000.10^3}{56}.\dfrac {178{,}29}{0{,}6} = 5306{,}25.10^3$ kJ \\ Ta có: $5306{,}25.10^3 = \dfrac {m.10^3.0{,}8}{12}.393{,}5 \Rightarrow m = 202{,}27$ kg.
\end{loigiaibth}
\def\writeANS{}
\begin{loigiaibth}{23}
  Gọi $\begin {cases} n_{C_3H_8} = a \, (mol) \\ n_{C_4H_{10}} = b \, (mol) \end {cases} \longrightarrow 44a + 58b = 12 \,\, (1)$ \\ $12$ gam X tỏa ra lượng nhiệt là: $2220.a + 2874.b = 597{,}6$ kJ $(2)$ \\ Từ $(1)$ và $(2)$ ta có hệ phương trình: $\begin {cases} 44a + 58b = 12 \\ 2220a + 2874b = 597{,}6 \end {cases} \longrightarrow \begin {cases} a = 0{,}075 \\ b = 0{,}15 \end {cases}$ \\ $\rightarrow $ a : b = 1 : 2.
\end{loigiaibth}
